\chapter{Introdução}
\label{cha:introducao}
\index{Introdução}

A calibração multivariada desempenha um papel fundamental em diversas áreas da análise química e industrial, permitindo a estimação de concentrações a partir de dados instrumentais multivariados. No contexto específico da análise de aço-ligas, a determinação precisa da composição elementar (Mn, Cr, Mo, Ni, Fe) é crucial para garantir a qualidade e as propriedades desejadas do material. No entanto, a complexidade dos dados multivariados provenientes de métodos instrumentais de análise exige o uso de técnicas de calibração multivariada para obter estimativas com erro de predição quantificável (RMSEP, MAE, R$^2$).

Este estudo investiga, em caráter exploratório, a aplicação de técnicas de reamostragem (\textit{bootstrap}), \textit{bagging} e \textit{subagging} na Regressão Linear Múltipla (RLM) para a calibração multivariada de amostras de aço-ligas. A pesquisa busca preencher a lacuna de conhecimento sobre a aplicação dessas técnicas de \textit{ensemble} nesse contexto específico, fornecendo subsídios empíricos exploratórios para a quimiometria e exemplificando a modernização de sistemas computacionais de apoio à análise de aço-ligas~\cite{SKOOG:1997,HARDLE:2007,TIMM:2002}.

Como otimizar a calibração multivariada de amostras de aço-ligas para obter previsões mais precisas dos elementos Mn, Cr, Mo, Ni e Fe, considerando a complexidade inerente à análise de dados multivariados, a necessidade de alta precisão e a exigência de uma arquitetura de \textit{software} modular e verificável por testes automatizados para sustentar essa análise?

Este estudo contribui em dois eixos: (a) a aplicação de \textit{bagging} e \textit{subagging} à calibração de aço-ligas; (b) a migração de um sistema legado verificada por testes automatizados. Essa estratégia evidencia, de forma reprodutível, como preservar a corretude de um núcleo algorítmico crítico durante a reescrita completa de uma arquitetura de \textit{software}.

A originalidade da pesquisa é dupla. No eixo quimiométrico, reside na aplicação de técnicas de \textit{bagging} e \textit{subagging}, ainda pouco exploradas na análise de aço-ligas.

No eixo de engenharia de \textit{software}, reside na condução de uma reescrita \textit{black-box} em três estágios (\textit{Desktop} Java, \textit{Web} Vaadin/Spring, Python/Streamlit), com arquitetura final desacoplada em três módulos de responsabilidade única e validação por 92 testes automatizados, que verificam a equivalência numérica do núcleo determinístico (RLM) e a consistência estatística dos métodos estocásticos entre o sistema legado e o sistema reescrito.

Dado o conjunto amostral reduzido — característica inerente à aquisição dos dados espectrométricos utilizados —, os resultados devem ser interpretados como evidências preliminares do comportamento dessas técnicas nesse domínio, e não como afirmações definitivas de superioridade geral de um método sobre outro.

Analisar e comparar o desempenho das técnicas de \textit{bagging} e \textit{subagging} aplicadas à Regressão Linear Múltipla (RLM) na calibração multivariada de dados de amostras de aço-ligas, investigando, em caráter exploratório, tanto o comportamento preditivo quanto a eficiência computacional de cada abordagem, a partir de uma arquitetura de \textit{software} desenvolvida de forma incremental e validada a cada etapa da migração.

Esta pesquisa tem como objetivos específicos:
\begin{description}
 \item[a)] Aplicar as técnicas de reamostragem \textit{bootstrap}, \textit{bagging} e \textit{subagging} na calibração de dados espectrométricos de amostras de aço-ligas, empregando a Regressão Linear Múltipla (RLM) como modelo base, visando a estruturação de um fluxo analítico reprodutível e padronizado~\cite{Canty2001,rathnayake2022}.

 \item[b)] Realizar uma comparação sistemática entre os algoritmos RLM clássica, RLM-\textit{bagging} e RLM-\textit{subagging}, por meio da avaliação criteriosa da acurácia preditiva e da eficiência computacional de cada abordagem na análise quantitativa dos elementos Mn, Cr, Mo, Ni e Fe.

 \item[c)] Reaproveitar, em Java (JavaSE-17), o sistema \textit{Desktop} de calibração multivariada com RLM-\textit{bagging} e RLM-\textit{subagging} — desenvolvido na disciplina de Programação Orientada a Objetos, com arquitetura monolítica (em um único arquivo), biblioteca \gls{EJML}\footnote{A biblioteca \textit{Efficient Java Matrix Library} (EJML) é uma biblioteca de álgebra linear de alto desempenho para Java, projetada para manipulação eficiente de matrizes densas e esparsas.} e interface gráfica Swing/AWT —, refinando suas estruturas de dados e tratando inconsistências nas métricas de avaliação do algoritmo para servir como referência de validação e base de código legada.

 \item[d)] Desenvolver e disponibilizar um protótipo \textit{Web} reescrito em Python 3.12 e Streamlit, adotando arquitetura modular desacoplada em três módulos de responsabilidade única (\texttt{app.py}, \texttt{components.py} e \texttt{models.py}), que automatize o fluxo analítico de calibração multivariada por RLM, RLM-\textit{bagging} e RLM-\textit{subagging}, permitindo o carregamento de matrizes de calibração e teste, a configuração de hiperparâmetros, a execução de validação cruzada 5-\textit{fold} por \textit{grid search} e a geração automática de relatórios com métricas preditivas (RMSEP, MAE, R$^2$) e gráficos interativos Plotly.

 \item[e)] Garantir a equivalência comportamental e a qualidade do \textit{software} reescrito por meio da especificação de requisitos formais e da implementação de uma suíte de 92 testes automatizados com \texttt{pytest}, abrangendo testes unitários, testes de caixa-preta baseados em saídas do sistema legado e testes de estresse com matrizes sintetizadas de alta dimensão.
\end{description}

Esse protocolo funciona como validador de equivalência entre o sistema legado e o sistema reescrito em dois níveis: equivalência numérica estrita para a RLM determinística e consistência estatística, dentro de tolerâncias explícitas, para o \textit{bagging} e o \textit{subagging}, cujas sequências pseudoaleatórias diferem entre as duas linguagens.

Este trabalho está organizado da seguinte forma. A Seção~\ref{cha:tecboot} apresenta a revisão da literatura sobre a Regressão Linear Múltipla (RLM), a técnica de reamostragem \textit{bootstrap} e os conceitos de distribuição empírica e princípio \textit{plug-in}, essenciais à compreensão das abordagens investigadas. A Seção~\ref{sec:teclarea} descreve as técnicas de \textit{bagging} e \textit{subagging}, incluindo fundamentos teóricos e algoritmos.